% !PRINTED 352
\weberpage{352}
%<<CONT
ist. Wir setzen auch voraus, dass die Curve vierter Ordnung frei von singulären Punkten sei\footnote{\gesperrt{Jacobi}, \glqq Ueber die Anzahl der Doppeltangenten ebener algebraischer Curven\grqq{}, \gesperrt{Crelle's} Journal, Bd.\,40 (1850). \gesperrt{Clebsch}, \glqq Bemerkung zu \gesperrt{Jacobi's} Beweis für die Anzahl der Doppeltangenten\grqq{}, ebend., Bd.\,63 (1864). Aus der ziemlich umfassenden Literatur über die Theorie der Doppeltangenten einer Curve vierter Ordnung sind die folgenden Schriften hervorzuheben. Zunächst mehrere Arbeiten von \gesperrt{Hesse} in \gesperrt{Crelle's} Journal, Bd.\,41, 49, 52. Ferner \gesperrt{Steiner}, Eigenschaften der Curven vierter Ordnung rücksichtlich ihrer Doppeltangenten, \gesperrt{Crelle's} Journal, Bd.\,49. \gesperrt{Aronhold}, Monatsber. d. Berl. Akademie 1864. \gesperrt{Geiser}, \glqq Ueber die Doppeltangenten einer ebenen Curve 4\,ten Grades\grqq{}, Math. Annalen, Bd.\,1 (1868). \gesperrt{Cayley}, \gesperrt{Crelle's} Journal, Bd.\,68 (1868). Auch das oben citirte Werk von \gesperrt{Salmon} ist hier hervorzuheben. In neuerer Zeit wurde die Theorie der Doppeltangenten im Zusammenhange mit der Theorie der \gesperrt{Abel'schen} Functionen weiter ausgebildet. \gesperrt{Riemann}, Gesammelte mathematische Werke (2.\,Aufl., 1892) XXXI, aus dem Nachlass. \gesperrt{Weber}, Theorie der \gesperrt{Abel'schen} Functionen vom Geschlecht 3, Berlin 1876. \gesperrt{Frobenius}, \gesperrt{Crelle's} Journal, Bd.\,99 u.\,103. Die algebraische Seite der Frage ist in zwei Abhandlungen: \gesperrt{Nöther}, Math. Annalen, Bd.\,15 und \gesperrt{Weber}, ebend., Bd.\,23 behandelt.}. Es möge jetzt $f(x_1, x_2, x_3)$ oder kürzer $f(x)$ eine ternäre Form 4\,ten Grades sein, die, gleich Null gesetzt, eine Curve vierter Ordnung ohne singulären Punkt darstellt, die wir die Curve $f$ nennen.

Setzen wir in dieser Gleichung
\weq{1}{f(x_1, x_2, x_3) = 0}
an Stelle der Variablen $x_1, x_2, x_3$ Ausdrücke von der Form

\[
x_1+t y_1, \quad x_2+t y_2, \quad x_3+t y_3,
\]

so ergiebt sich eine Gleichung 4\,ten Grades in Bezug auf $t$:
\weq{2}{f(x+ty) = 0,}
deren Wurzeln, wenn $(x)$ und $(y)$ zwei feste Punkte sind, in $x+ty$ eingesetzt, die Coordinaten der Schnittpunkte der Verbindungslinie der Punkte $(x), (y)$, die wir als die Linie $(x,y)$ bezeichnen wollen, mit der Curve $f$ geben.

Es werde nun die Function $f(x+ty)$ nach Potenzen von $t$ geordnet. Dies giebt (nach Bd.\ I, \S.\,60):
\weq{3}{f(x+yt) = P_0(x,y)+4tP_1(x,y)+6t^2P_2(x,y)+4t^3P_3(x,y)+t^4P_4(x,y),}
worin die $P_\nu(x,y)$ die Polaren von $f(x)$ sind, also homogene
%CONT>>
